\documentclass{article}
\usepackage[left=2cm, right=2cm, top=2cm, bottom=2cm]{geometry}
\usepackage{graphicx}
\usepackage{amsmath}
\usepackage{amssymb}
\usepackage{amsthm}
\usepackage{fancyhdr}
\usepackage{verbatim}
\usepackage{listings}
\usepackage{xcolor}
\usepackage{pdfpages}
\usepackage{cancel}
\usepackage{physics}
\usepackage{esint}
\usepackage{braket}

\lstset{
    language=C++,
    basicstyle=\ttfamily\footnotesize,
    keywordstyle=\color{blue},
    commentstyle=\color{green},
    stringstyle=\color{red},
    numbers=left,
    numberstyle=\tiny\color{gray},
    frame=single,
    breaklines=true,
    captionpos=b,
}


\title{MTH 446: Lecture \# 19}
\author{Cliff Sun}

\newtheorem{theorem}{Theorem}[section]
\newtheorem{lemma}[theorem]{Lemma}
\newtheorem{definition}[theorem]{Definition}
\newtheorem{conjecture}[theorem]{Conjecture}
\newtheorem{proposition}[theorem]{Proposition}
\newtheorem{corollary}[theorem]{Corollary}
\newtheorem{one minute paper}[theorem]{One Minute Paper}

\pagestyle{fancy}
\lhead{\textbf{\thepage}\ \ \nouppercase{\rightmark}}
\chead{MTH 446: Lecture \# 19}
\rhead{Cliff Sun}

\begin{document}

\maketitle

\section*{Lecture Span}
\begin{enumerate}
    \item Riemann Integration
\end{enumerate}

\section*{Integration}

\begin{equation}
    \int_a^b w dt = \int_a^b u + i\int_a^b v
\end{equation}

Therefore, 

\begin{equation}
    \Re \int w dt = \int Re w
\end{equation}

\begin{equation}
    \Im \int w dt = \Im Re w
\end{equation}

Consider a vector function in $\mathbb{R}^2$, then 

\begin{equation}
    \int_a^b f = (\int_a^b x, \int_a^b y)
\end{equation}

Three properties:
\begin{enumerate}
    \item $\int_a^b (f + g) = \int f + \int g$
    \item $\Re \int w dt = \int Re w$
    \item $\Im \int w dt = \Im Re w$
    \item $\int z_0 f = z_0 \int f$
\end{enumerate}

We consider piecewise continuous functions. 
A function $h$ is piecewise continuous if $h$ is continuous at every point of $[a,b]$ except for a finite set of points $\{x_0,\dots, x_n\}$. 

Here, $x_0 = a$ and $x_n = b$ WLOG. $h$ also has a one-sided limit. Suppose $x_j \in \{a,b\}$, then the left and right limit of $h$ as $x \to x_j$ must exist.

It is known that every piecewise continuous function is Riemann integrable. The fundamental theorem of calculus: 

\begin{equation}
    \int_a^b w' = w(b) - w(a)
\end{equation}

This is true even for complex functions. Now, let $w$ be piecewise continuous. Then, $|w|$ is also piecewise continuous. 

\begin{theorem}
    Let $w:[a,b] \to \mathbb{C}$ be a piecewise continuous function. Then the following inequality holds:
    \begin{equation}
        |\int_a^b w dt | \leq \int_a^b |w| dt
    \end{equation}
\end{theorem}

\begin{proof}
    If $\int_a^bwdt = 0$, then we are done. Next, assume $\int_a^b w dt \neq 0$. Then, 
    \begin{align}
        \int_a^b w dt &= r_0 \exp(i\theta_0)\\
        |\int_a^b w dt| &= r_0 \\
        &= \Re\left( \int_a^b \exp(-i\theta_0)wdt \right)\\
        &= |\int_a^b \Re(\exp(-i\theta_0)w dt) \\
        &\leq \int_a^b |\Re(\exp(-i\theta_0)w)|dt \\
        |\int_a^b w dt| = r_0 &\leq \int_a^b |\Re(\exp(-i\theta_0)w)|dt \leq \int_a^b |\exp(-i\theta_0)w|dt = \int_a^b |w| dt
    \end{align}
\end{proof}

\subsection*{Example}

Let 

\begin{equation}
    P_n(x) = \frac{1}{\pi}\int_0^\pi (x + i\sqrt{1 - x^2}\cos\theta)^n d\theta
\end{equation}

Then prove that $|P_n| \leq 1$. Then 

\begin{equation}
    |P_n| \leq \frac{1}{\pi}\int_0^\pi|x + i\sqrt{1-x^2}\cos\theta|^n d\theta
\end{equation}

Moreover, 

\begin{equation}
    |x + i\sqrt{1 - x^2}\cos\theta|^2 = x^2 + (1 - x^2)\cos^2\theta \leq x^2 + 1 - x^2 = 1
\end{equation}

Therefore, 

\begin{equation}
    |x + i \sqrt{1 - x^2}\cos\theta|^n \leq 1
\end{equation}

Thus, 

\begin{equation}
    |P_n| \leq \frac{1}{\pi}\int_0^\pi d\theta = 1
\end{equation}

\begin{definition}
    A \underline{Curve} is a continuous function $\gamma:[a,b] \to \mathbb{C}$
\end{definition}

\begin{definition}
     A curve is \underline{closed} if $\gamma(a) = \gamma(b)$
\end{definition}

\begin{definition}
    A curve $\gamma$ is smooth if 
    \begin{enumerate}
        \item $\gamma \in C^1$
        \item $\gamma' \neq 0$ for all $t$.  
    \end{enumerate}
\end{definition}

Example, $\gamma = (t^2, t^3)$ is not smooth. 

\begin{definition}
    A curve $g$ is called a \underline{simple curve} if $t_1,t_2 \in [a,b]$ and $t_1 \neq t_2$, then $\gamma(t_1) \neq \gamma(t_2)$. 
\end{definition}

You can have a simple, closed curve. The arc length of a curve is 

\begin{equation}
    s(\gamma) = \int_a^b \sqrt{x'^2 + y'^2}dt
\end{equation}

Note that $|\gamma'| = \sqrt{x'^2 + y'^2}$, and therefore, in complex notation, we have that:

\begin{equation}
    s(\gamma) = \int_a^b |g'| dt
\end{equation}



\end{document}